\title{Controlling Text-to-Image Diffusion by Orthogonal Finetuning}

\begin{abstract}
State-of-the-art text-to-image diffusion models have demonstrated remarkable success in synthesizing photorealistic visuals based on textual descriptions. However, developing effective mechanisms to direct or modulate these models for a range of specific tasks remains a challenging and unsolved issue. In this work, we propose a theoretically grounded adaptation strategy—Orthogonal Finetuning~(OFT)—to tailor text-to-image diffusion models for diverse downstream applications. In contrast to current approaches, OFT can be shown to maintain hyperspherical energy, encapsulating the pairwise associations among neurons situated on the unit hypersphere. Our findings indicate that preserving this hyperspherical structure is essential for maintaining the semantic generative capacities of text-to-image diffusion models. To further enhance finetuning robustness, we introduce Constrained Orthogonal Finetuning (COFT), which incorporates a supplementary constraint on the hypersphere's radius. Our focus encompasses two principal text-to-image adaptation scenarios: subject-driven generation, aiming to render subject-specific outputs from a limited set of subject images combined with textual guidance, and controllable generation, in which the model must also respond to explicit control signals. Experimental results establish that the OFT framework yields superior generation quality and faster convergence when compared to previous finetuning strategies.
\end{abstract}

\section{Introduction}

Cutting-edge text-to-image diffusion models~\cite{saharia2022photorealistic,ramesh2022hierarchical,rombach2022high} have set new benchmarks in generating high-quality images steered by text instructions. Despite these advancements, textual prompts may lack the necessary precision and can fall short in offering detailed and exacting control over the resulting images. To confront this issue, our study centers on two prominent challenges in text-to-image generation:

\begin{itemize}[leftmargin=*,nosep]

    \item \textbf{Subject-driven generation}~\cite{ruiz2023dreambooth}: Provided with only a handful of images depicting a particular subject, the objective is to synthesize new images portraying that subject within varied contexts, all guided by an accompanying text prompt.
    \item \textbf{Controllable generation}~\cite{zhang2023adding,mou2023t2i}: When given supplemental control data (\emph{e.g.}, canny edge outlines, segmentation masks), the model is tasked with creating images that comply with both the specified control information and a descriptive prompt.
\end{itemize}

Both tasks are fundamentally concerned with optimizing the finetuning process of text-to-image diffusion models in such a way that the original generative quality acquired through pretraining is retained. The key objectives for a robust finetuning approach can be summarized as: (1) \emph{training efficiency}, which encompasses minimizing both the count of parameters updated and the number of training epochs required; and (2) \emph{generalizability preservation}, which refers to maintaining the model's capacity to generate high-quality and diverse outputs. Conventional finetuning methods generally adopt one of two strategies: updating model weights incrementally with a small learning rate (\emph{e.g.}, \cite{ruiz2023dreambooth}) or introducing a lightweight, re-parameterized module to the model (\emph{e.g.}, \cite{hulora2022,zhang2023adding}). However, these approaches lack formal guarantees for preserving the generative capabilities established during pretraining. Furthermore, there remains an absence of rigorous theoretical foundations guiding the selection of effective finetuning procedures or the appropriate setting of hyperparameters such as the number of training epochs. One of the core challenges arises from not having a clear metric to assess how much of the original generative power is maintained after finetuning. Current approaches often operate under the assumption that a smaller Euclidean distance between the updated and original models correlates with better retention of pretrained abilities, leading to a preference for low learning rates. Although this relationship may sometimes be valid, we contend that Euclidean distance alone is insufficient to capture semantic preservation comprehensively. Therefore, introducing a more structurally informed metric for evaluating the divergence between finetuned and pretrained models has the potential to significantly enhance both the retention of generative performance and the reliability of the finetuning process.

Drawing on empirical findings that hyperspherical similarity effectively captures semantic content~\cite{liu2017deep,liu2018decoupled,chen2020angular}, we leverage hyperspherical energy~\cite{liu2018learning} to represent the pairwise relationships among neurons. This metric, calculated as the aggregate hyperspherical similarity (such as cosine similarity) across all pairs of neurons within a single layer, reflects the degree of uniform distribution of neurons on the unit hypersphere~\cite{liu2021learning}. Our underlying assumption is that an optimally finetuned model should maintain minimal deviation in hyperspherical energy from its pretrained counterpart. An intuitive strategy involves incorporating a regularization term that enforces constancy of hyperspherical energy during finetuning. However, this method does not ensure that the energy discrepancy will be effectively reduced. 

Instead, we exploit a key invariance property of hyperspherical energy: the pairwise hyperspherical similarity remains unchanged under a common orthogonal transformation applied to all neurons. Building upon this property, we introduce Orthogonal Finetuning~(OFT), a method designed to adapt extensive text-to-image diffusion models to new tasks while keeping their hyperspherical energy unaltered. The core concept involves learning a layer-wide orthogonal transformation for neuron representations, thereby preserving pairwise angular relationships. This approach can alternatively be seen as transforming the canonical basis for neuron vectors within the same layer. Additionally, acknowledging that a reduced Euclidean distance between the pretrained and finetuned models is indicative of better retention of original model performance, we present a refined version called Constrained Orthogonal Finetuning~(COFT), which restricts the updated model to lie within a hypersphere of constant radius centered around the pretrained neuron states.

The rationale for employing orthogonal transformations in neuron finetuning is partially motivated by concepts from the 2D Fourier transform, where an image can be represented through its magnitude and phase spectra. Importantly, the phase spectrum—corresponding to the angular relation between the input and the basis—encodes most of the semantic content. For instance, when reconstructing an image, retaining its phase spectrum while substituting the magnitude spectrum with a random one can still yield an image that largely maintains the original semantics. This observation implies that altering the orientations of neurons is central to effecting meaningful semantic modifications in generated images, which is a fundamental objective in both subject-driven and controlled image synthesis. Nevertheless, if the neuron directions are modified with excessive flexibility, there is a significant risk of compromising the generative abilities established during pretraining. To mitigate this, we advocate preserving the pairwise angles among all neurons, grounded in the hypothesis that these angular relationships are essential for maintaining the representational capacity of neural networks. Following this line of reasoning, it is thus logical to learn a layer-wise, shared orthogonal transformation that maintains the hyperspherical energy across neurons within each layer, thereby conserving the geometrical properties of the pretrained model.

Our approach is further inspired by orthogonal over-parameterized training~\cite{liu2021orthogonal}, a method which initializes classification neural networks randomly and employs orthogonal transformations during training. This strategy takes advantage of the fact that random initializations typically result in low expected hyperspherical energy, and the objective in~\cite{liu2021orthogonal} is to maintain this low energy throughout training (as lower energy correlates with improved classification generalization~\cite{liu2018learning,lin2020regularizing}). Their findings indicate that orthogonal transformations possess sufficient flexibility to enable the training of neural networks capable of generalizing well on classification tasks. By contrast, our work targets the finetuning of text-to-image diffusion models to enhance both controllability and downstream generative capacity. There are two primary distinctions between OFT and~\cite{liu2021orthogonal}: firstly, while~\cite{liu2021orthogonal} is concerned with minimizing hyperspherical energy, OFT is designed to conserve the hyperspherical energy characteristic of the pretrained network, thereby safeguarding its structural integrity during finetuning. In the context of diffusion models, further minimizing this energy may disrupt inherent semantic relationships. Secondly, OFT explicitly strives to limit its deviation from the original pretrained configuration, introducing a constrained variant to achieve this; in contrast,~\cite{liu2021orthogonal} does not incorporate such a constraint. Balancing adaptability with stability is crucial in the finetuning process, and we contend that the OFT paradigm realizes this balance effectively. Our main contributions are outlined as follows:

\begin{itemize}[leftmargin=*,nosep]

    \item We introduce a new finetuning technique—Orthogonal Finetuning—for enhancing the controllability of text-to-image diffusion models. To address stability concerns, we also present a constrained version that restricts the angular shift relative to the original pretrained parameters.
    \item In contrast with established finetuning approaches, OFT modifies the model while guaranteeing the preservation of hyperspherical energy—a property that our experiments show to be vital for maintaining the generative semantics encoded in the pretrained model.
    \item We evaluate OFT on two application areas: subject-driven generation and controllable generation. Through extensive experiments, we observe notable gains over previous methods in terms of output quality, training efficiency, and the robustness of finetuning. OFT also demonstrates improved sample efficiency, requiring fewer images and training epochs to achieve convergence.
    \item For the controllable generation task, we propose a novel control consistency metric to assess the degree of controllability. This metric operates by inferring the control signal from the output image and comparing it to the original control prompt, providing quantitative evidence for OFT's superior controllability.
\end{itemize}

\section{Related Work}

\textbf{Text-to-image diffusion models}. Text-to-image synthesis has seen significant advancements~\cite{saharia2022photorealistic,ramesh2022hierarchical,rombach2022high,gu2022vector,nichol2022glide}, owing to breakthroughs in diffusion-based generation~\cite{ho2020denoising,song2021score,song2021denoising,dhariwal2021diffusion} and progress in learning joint vision-language representations~\cite{su2020vlbert,lu2019vilbert,radford2021learning,wang2021simvlm,li2022blip,li2023blip,alayrac2022flamingo,schuhmann2022laion}. For instance, GLIDE~\cite{nichol2022glide} and Imagen~\cite{saharia2022photorealistic} train diffusion models directly in pixel space. While GLIDE optimizes the text encoder jointly with a diffusion prior leveraging paired image-text data, Imagen relies on a pretrained, frozen text encoder. Latent-space diffusion modeling is used by Stable Diffusion~\cite{rombach2022high} and DALL-E2~\cite{ramesh2022hierarchical}; the former utilizes VQ-GAN~\cite{esser2021taming} to construct a learned visual latent space, and the latter incorporates CLIP~\cite{radford2021learning} for a unified latent embedding of image and text. Beyond diffusion-based architectures, text-to-image generation has also been explored using generative adversarial networks~\cite{reed2016generative,zhang2017stackgan,xu2018attngan,li2019controllable} and autoregressive strategies~\cite{ramesh2021zero,ding2021cogview,wu2022nuwa,yuscaling}. Notably, OFT is designed to be agnostic to specific model architectures and thus is applicable across different text-to-image diffusion models.

\textbf{Subject-driven generation}. To restrict alterations to the target subject, works such as \cite{avrahami2022blended,nichol2022glide} introduce user-provided masks as supplementary conditions. Inversion-based frameworks~\cite{choi2021ilvr,dhariwal2021diffusion,ramesh2022hierarchical,gal2022image} enable context manipulation while keeping the main subject unchanged, and \cite{hertz2022prompt} is capable of executing both localized and global edits without requiring any user-supplied mask. Nevertheless, these techniques often fail to capture fine identity-specific attributes of the subject. The Pivotal Tuning method~\cite{roich2022pivotal} adapts the generator around an inverted latent point while incorporating additional regularization to safeguard subject identity. A related direction, explored by \cite{nitzan2022mystyle}, involves learning a custom generative face prior from a dataset of one person's faces. In \cite{casanova2021instance}, a method is presented to generate diverse variants of an instance, albeit sometimes at the cost of detailed instance features. DreamBooth~\cite{ruiz2023dreambooth} employs a personalized token and a small set of subject images to finetune a diffusion model, combining reconstruction and class-specific prior preservation losses. Building on DreamBooth, OFT introduces orthogonal transformations during finetuning, departing from standard gradient-based updates with small learning rates.

\textbf{Controllable generation}. Image-to-image translation is commonly considered a variant of controllable generation. Traditionally, this has been approached using conditional generative adversarial networks~\cite{isola2017image,zhu2017toward,wang2018high,park2019semantic,choi2018stargan}. Recent developments have seen diffusion models applied to image-to-image translation tasks~\cite{saharia2022palette,voynov2022sketch,wang2022pretraining}. Among newer advances, ControlNet~\cite{zhang2023adding} enhances a pretrained diffusion model by fine-tuning it to accept extra conditioning signals, leading to substantial improvements in controlled generation capabilities. T2I-Adapter~\cite{mou2023t2i}, a related contemporary method, also finetunes pretrained diffusion models for increased control over outputs. Building on the protocols established by~\cite{zhang2023adding,mou2023t2i}, our work introduces OFT for finetuning pretrained diffusion models, achieving superior controllability with reduced training data requirements and a smaller number of finetuned parameters. Notably, OFT incurs no extra computation during inference at test time.

\textbf{Model finetuning}. The practice of refining large, pretrained models for specific downstream tasks is rapidly gaining traction~\cite{devlin2018bert,brown2020language,he2022masked}. A variety of adaptation techniques, including those outlined in~\cite{hulora2022,houlsby2019parameter,pfeiffer2020adapterfusion}, are extensively researched in natural language processing as instances of finetuning. LoRA~\cite{hulora2022}, closely related to OFT, utilizes a low-rank approximation for additive weight adjustments during model adaptation. By contrast, OFT updates neuron weights through layer-wise shared orthogonal transformations in a multiplicative way. This mechanism rigorously maintains the angles between neurons within the same layer, providing improved training stability.

\section{Orthogonal Finetuning}

\subsection{Why Does Orthogonal Transformation Make Sense?}

We begin by examining the rationale for employing orthogonal transformation in the finetuning of text-to-image diffusion models. To address this, we break the issue into two focal questions: (1) what motivates the adjustment of neuron angles (i.e., their directions) during finetuning, and (2) the reasoning behind choosing orthogonal transformation as the means for modifying these angles.

Addressing the first question, we take cues from prior empirical findings reported in \cite{liu2018decoupled,chen2020angular}, which reveal that differences in feature angles can effectively capture the semantic disparity between representations. SphereNet~\cite{liu2017deep} further establishes that constraining neural network weights to reside on a unit hypersphere retains comparable representational power and enhances generalization, indicating that the directional components of neurons predominantly encode data semantics. To further illustrate the significance of neuron angles, we design a simple experiment depicted in Figure~\ref{fig:toy_ae}, wherein a conventional convolutional autoencoder is trained on a dataset of flower images. During training, we adopt the usual inner product for feature map computation ($z$ denotes the output of the convolution filter $\bm{w}$ when applied to input $\bm{x}$ in the sliding window). In the evaluation phase, we generate the feature map using three approaches: (a) the same inner product as in training, (b) leveraging only magnitude information, and (c) relying solely on angular information. Results in Figure~\ref{fig:toy_ae} reveal that neuron angles are nearly sufficient to reconstruct input images with high fidelity, while magnitude alone contributes negligible information. It is important to note that the cosine activation (c) is not used during training; all learning is based on the inner product. These findings suggest that neuron directions are primarily responsible for retaining the semantic content of input images. Consequently, adjusting neuron directions appears to be a promising strategy for manipulating the semantics encoded in neural representations.

Regarding the second question, a direct method to update neuron directions involves jointly rotating or reflecting all neurons within a layer, thereby introducing an orthogonal transformation. Although more intricate angular transformations—where neurons are individually rotated by distinct angles—could offer increased adaptability, orthogonal transformations strike an effective balance between versatility and structural regularity. Additionally, \cite{liu2021orthogonal} demonstrates the expressive capacity of orthogonal transformations for neural network optimization. To substantiate this claim, we conduct a regularization experiment modeled after the ControlNet~\cite{zhang2023adding} setup, with outcomes presented in Figure~\ref{fig:ortho}. The results show that applying our standard OFT method yields stable performance and precise control by the end of training (epoch 20). In contrast, removing the orthogonality constraint from OFT results in a failure to generate realistic images or achieve control. This experiment underscores the critical role of orthogonality in ensuring the effectiveness of OFT.

\subsection{General Framework}

Traditional finetuning methods generally employ gradient descent with a reduced learning rate to adjust either the entire model or select layers. By adopting a low learning rate, these methods implicitly limit how much the finetuned model strays from the original pretrained weights, effectively minimizing $\thickmuskip=2mu \medmuskip=2mu \| \bm{M} - \bm{M}^0\|$—where $\bm{M}$ and $\bm{M}^0$ represent the finetuned and pretrained parameters, respectively. While this serves as a natural regularization, it may still grant excessive freedom, particularly for large-scale models. To counteract this, LoRA augments this procedure by enforcing a low-rank constraint on the update: $\thickmuskip=2mu \medmuskip=2mu \text{rank}(\bm{M} - \bm{M}^0)=r'$, with $r'$ chosen to be small. In contrast, OFT imposes a different kind of restriction, specifically on pairwise neuron relationships via hyperspherical energy: $\thickmuskip=2mu \medmuskip=2mu \| \text{HE}(\bm{M})-\text{HE}(\bm{M}^0) \|=0$, where $\text{HE}(\cdot)$ quantifies the hyperspherical energy associated with a weight matrix. 

To provide further intuition, consider a fully connected layer characterized by $\thickmuskip=2mu \medmuskip=2mu \bm{W}=\{\bm{w}_1,\cdots,\bm{w}_n\}\in\mathbb{R}^{d\times n}$, where each column $\bm{w}_i\in\mathbb{R}^{d}$ denotes the $i$-th neuron, and $\bm{W}^0$ refers to the pretrained weights. The output, denoted as $\thickmuskip=2mu \medmuskip=2mu \bm{z}\in\mathbb{R}^n$, is obtained through $\thickmuskip=2mu \medmuskip=2mu \bm{z}=\bm{W}^\top\bm{x}$ for input $\thickmuskip=2mu \medmuskip=2mu \bm{x}\in\mathbb{R}^d$. Within this framework, OFT can be understood as attempting to reduce the hyperspherical energy gap between the modified and original weights:

\begin{equation}\label{eq:oft_principle}
\footnotesize
\min_{\bm{W}} \left\| \text{HE}(\bm{W})-\text{HE}(\bm{W}^0)\right\|~~\Leftrightarrow~~\min_{\bm{W}} \bigg{\|} \sum_{i\neq j} \|\hat{\bm{w}}_i-\hat{\bm{w}}_j\|^{-1}-\sum_{i\neq j} \|\hat{\bm{w}}^0_i-\hat{\bm{w}}^0_j\|^{-1}\bigg{\|}
\end{equation}

Here, each normalized neuron is written as $\thickmuskip=2mu \medmuskip=2mu \hat{\bm{w}}_i:=\bm{w}_i/\|\bm{w}_i\|$, and hyperspherical energy for a layer $\bm{W}$ is defined as $\thickmuskip=2mu \medmuskip=2mu \text{HE}(\bm{W}):= \sum_{i\neq j} \|\hat{\bm{w}}_i-\hat{\bm{w}}_j\|^{-1}$. Notably, the global optimum for Eq.~\eqref{eq:oft_principle} is zero, achievable whenever $\bm{W}$ and $\bm{W}^0$ are connected by a rotation or reflection; in other words, if $\thickmuskip=2mu \medmuskip=2mu \bm{W}=\bm{R}\bm{W}^0$ for some orthogonal matrix $\thickmuskip=2mu \medmuskip=2mu \bm{R}\in\mathbb{R}^{d\times d}$ (where determinant $1$ or $-1$ corresponds to rotation or reflection, respectively). This insight underpins OFT, which restricts finetuning to optimizing such layer-wise orthogonal transformations. Specifically, OFT seeks an orthogonal matrix $\thick

\begin{equation}\label{eq:oft_general}
    \footnotesize
    \bm{z}=\bm{W}^\top\bm{x}=(\bm{R}\cdot\bm{W}^0)^\top\bm{x},~~~\text{s.t.}~\bm{R}^\top\bm{R}=\bm{R}\bm{R}^\top=\bm{I}
\end{equation}

Here, $\bm{W}$ stands for the weight matrix after OFT finetuning, and $\bm{I}$ represents the identity matrix. An illustration of OFT appears in Figure~\ref{fig:block}. Analogous to the zero initialization used in LoRA, it is necessary for OFT to update the pretrained model starting from $\bm{W}^0$. To realize this, the orthogonal matrix $\bm{R}$ is initialized as the identity, ensuring that the finetuned model begins from the original pretrained parameters. To maintain the orthogonality constraint for $\bm{R}$ during optimization, differentiable orthogonalization approaches, such as those described in \cite{liu2021orthogonal,lezcano2019cheap}, can be employed. We will elaborate on methods for enforcing orthogonality efficiently.

\subsection{Efficient Orthogonal Parameterization}\label{sect:eff_ortho}

Traditional orthogonalization techniques, like the Gram-Schmidt process, are differentiable but generally not computationally viable for large-scale applications~\cite{liu2021orthogonal}. To achieve better computational efficiency, Cayley parameterization is utilized to construct the orthogonal matrix. Specifically, we define the orthogonal matrix by setting $\thickmuskip=2mu \medmuskip=2mu \bm{R}=(\bm{I}+\bm{Q})(I-\bm{Q})^{-1}$, where the matrix $\bm{Q}$ is skew-symmetric and fulfills $\thickmuskip=2mu \medmuskip=2mu \bm{Q}=-\bm{Q}^\top$. This method is computationally efficient but imposes a mild constraint: Cayley parameterization only generates orthogonal matrices with determinant $1$—that is, elements of the special orthogonal group. Empirically, we observe that this restriction does not impact practical performance. Nevertheless, when the dimension $d$ is large, parameterizing $\bm{R}$ with the full Cayley transform can still be parameter-inefficient. To reduce parameter costs, we represent $\bm{R}$ as a block-diagonal matrix partitioned into $r$ blocks, which takes the form:

\begin{equation}
    \footnotesize
    \bm{R}=\text{diag}(\bm{R}_1,\bm{R}_2,\cdots,\bm{R}_r)=\begin{bmatrix}
    \bm{R}_1\in \text{O}(\frac{d}{r}) &  &\\
    &\ddots & \\
    & & \bm{R}_r\in \text{O}(\frac{d}{r})
    \end{bmatrix}\in \text{O}(d)
\end{equation}

Here, $\text{O}(d)$ represents the group of orthogonal matrices of size $d$, with $\thickmuskip=2mu \medmuskip=2mu \bm{R}\in\mathbb{R}^{d\times d}$ and each block $\thickmuskip=2mu \medmuskip=2mu \bm{R}_i\in\mathbb{R}^{d/r\times d/r}$ for all $i$ also being orthogonal. In the special case where $\thickmuskip=2mu \medmuskip=2mu r=1$, the block-diagonal orthogonal matrix is equivalent to a standard, unconstrained orthogonal matrix. For an orthogonal matrix of size $\thickmuskip=2mu \medmuskip=2mu d\times d$, the count of independent parameters equals $\thickmuskip=2mu \medmuskip=2mu d(d-1)/2$, which leads to $\mathcal{O}(d^2)$ complexity. By employing an $r$-block diagonal structure, the total parameter count becomes $\thickmuskip=2mu \medmuskip=2mu d(d/r-1)/2$, with complexity reduced to $\thickmuskip=2mu \medmuskip=2mu \mathcal{O}(d^2/r)$. Further parameter reduction is possible if we opt to share blocks, that is, if we set $\thickmuskip=2mu \medmuskip=2mu\bm{R}_i=\bm{R}_j$ for all $i\neq j$, shrinking the parameter complexity to $\mathcal{O}(d^2/r^2)$. Regardless of these efficiency strategies, the resulting $\bm{R}$ continues to satisfy the orthogonality condition, thereby ensuring that the hyperspherical energy is retained.

We next compare the parameter efficiency between OFT and LoRA. In the case of LoRA, when a low-rank parameter $r'$ is used, the trainable parameter count is given by $\thickmuskip=2mu \medmuskip=2mu r'(d+n)$. Suppose both $r$ and $r'$ scale with the model dimension $d$, such as taking $\thickmuskip=2mu \medmuskip=2mu r=r'=\alpha d$, where $0<\alpha\leq 1$ is constant; then, LoRA's parameter complexity is $\mathcal{O}(d^2+dn)$, while that of OFT becomes $\mathcal{O}(d)$. To clarify the distinction, consider a weight matrix of shape $\thickmuskip=2mu \medmuskip=2mu 128\times 128$; if $\thickmuskip=2mu \medmuskip=2mu r'=8$ for LoRA, this approach requires $2,048$ parameters, whereas OFT, with $\thickmuskip=2mu \medmuskip=2mu r=8$ (and without block sharing), only needs $960$ trainable parameters.

\subsection{Constrained Orthogonal Finetuning}

The flexibility of OFT can be further restricted by ensuring that finetuning does not cause the model to stray far from the initial pretrained solution. This more conservative approach, called COFT, employs the following computation:

\begin{equation}\label{eq:coft_general}
    \footnotesize
    \bm{z}=\bm{W}^\top\bm{x}=(\bm{R}\cdot\bm{W}^0)^\top\bm{x},~~~\text{s.t.}~\bm{R}^\top\bm{R}=\bm{R}\bm{R}^\top=\bm{I},~\norm{\bm{R}-\bm{I}}\leq \epsilon
\end{equation}

This formulation simultaneously enforces orthogonality and confines $\bm{R}$ to be within an $\epsilon$-ball of the identity matrix. Orthogonality is handled through Cayley parameterization (Section~\ref{sect:eff_ortho}), though incorporating the $\epsilon$-neighborhood constraint is less straightforward with Cayley-parameterized matrices. To better understand the effect of the Cayley transform, the Neumann series is used to approximate $\thickmuskip=2mu \medmuskip=2mu \bm{R}=(\bm{I}+\bm{Q})(I-\bm{Q})^{-1}$, yielding the expression $\thickmuskip=2mu \

series converges in the operator norm), which permits us to shift the constraint $\thickmuskip=2mu \medmuskip=2mu\|\bm{R}-\bm{I}\|\leq\epsilon$ into the Cayley transform framework. The corresponding equivalent formulation becomes $\thickmuskip=2mu \medmuskip=2mu \|\bm{Q}-\bm{0}\|\leq\epsilon'$, with $\bm{0}$ signifying the zero matrix and $\epsilon'$ being a separate tolerance parameter distinct from $\epsilon$. This new condition on $\bm{Q}$ is straightforward to impose using projected gradient descent methods. For initialization, setting $\bm{Q}$ as the zero matrix ensures that the orthogonal matrix $\bm{R}$ begins as the identity. COFT is thus interpretable as the superposition of two explicit constraints: maintaining minimal hyperspherical energy change and bounding deviation from the base (pretrained) model. The latter is usually an implicit assumption in standard finetuning strategies, but COFT formalizes it as an explicit requirement. While OFT demonstrates strong performance, our observations suggest that the explicit deviation constraint in COFT provides even greater finetuning stability than OFT alone. As illustrated in Figure~\ref{fig:eps_coft}, the choice of $\epsilon$ modulates COFT’s performance. A larger $\epsilon$ yields finetuned models akin to those from OFT; reducing $\epsilon$ drives the finetuned model's behavior closer to the original pretrained text-to-image diffusion model.

\subsection{Re-scaled Orthogonal Finetuning}

To build on the standard OFT, we introduce an additional learnable scaling factor per neuron, motivated by the observation that rescaling neurons leaves the hyperspherical energy invariant (since magnitudes are normalized). The feedforward computation thus becomes: $\thickmuskip=2mu \medmuskip=2mu\bm{z}=(\bm{R}\bm{W}^0\bm{D})^\top\bm{x}$\footnote{\textbf{Errata}: The forward operation for re-scaled OFT was erroneously reported as $\thickmuskip=2mu \medmuskip=2mu\bm{z}=(\bm{D}\bm{R}\bm{W}^0)^\top\bm{x}$ in the NeurIPS camera ready version, but the underlying implementation is correct and thus the outcomes in Appendix~\ref{app:rescaled} remain valid.} where $\thickmuskip=2mu \medmuskip=2mu \bm{D}=\text{diag}(s_1,\cdots,s_n)\in\mathbb{R}^{n\times n}$ represents a trainable diagonal matrix with strictly positive entries $s_1,\ldots,s_n$. Unlike OFT as defined in Eq.~\eqref{eq:oft_general}, which involves optimization over only $\bm{R}$, the re-scaled variant permits both $\bm{R}$ and $\bm{D}$ to adapt during learning. This extension augments OFT’s expressivity, yet introduces a negligible increase in parameter count. In our main empirical study, we retain the classic OFT to highlight the power of pure orthogonal transformation, but observe that the re-scaled approach frequently attains even better outcomes (see Appendix~\ref{app:rescaled} for supporting results).

\section{Intriguing Insights and Discussions}

\textbf{OFT is not specific to a particular neural architecture}. In principle, OFT can be applied to any neural network structure. For instance, within Transformer models, LoRA is commonly deployed to adapt attention matrices~\cite{hulora2022}. To ensure a fair comparison with LoRA, our experiments use OFT to modify only the attention parameters. Moreover, OFT is highly suitable for finetuning convolutional layers, as the block-diagonal nature of $\bm{R}$ provides interesting properties specific to convolutions—something LoRA lacks. If the number of blocks matches the number of input channels, each block operates independently on a single neuron channel, analogous to depth-wise convolutional kernels~\cite{chollet2017xception}. Alternatively, sharing blocks across all channels allows OFT to perform an orthogonal transformation identically across neuron channels. For empirical findings on using OFT with convolutional layers, see our exploratory analysis in Appendix~\ref{app:convolution}.

\textbf{Connection to LoRA}. LoRA introduces a low-rank matrix to the model, which helps preserve the information encoded in the pretrained weight matrix by avoiding substantial shifts in those weights. By contrast, OFT ensures that the transformation applied to the pretrained weight matrix remains orthogonal (thus full-rank), thereby safeguarding the integrity of the pretrained knowledge. The OFT forward computation can be reformulated as $\thickmuskip=2mu \medmuskip=2mu \bm{z}=(\bm{R}\bm{W}^0)^\top\bm{x}=(\bm{W}^0+(\bm{R}-\bm{I})\bm{W}^0)^\top\bm{x}$, in which $\thickmuskip=2mu \medmuskip=2mu (\bm{R}-\bm{I})\bm{W}^0$ acts in a similar capacity to LoRA’s low-rank weight modification. Since $\bm{W}^0$ is ordinarily full-rank, when $\thickmuskip=2mu \medmuskip=2mu \bm{R}-\bm{I}$ has low rank, OFT similarly introduces a low-rank update to the weights. Analogous to the rank parameter $r'$ used in LoRA for controlling trainable parameters, OFT employs a diagonal block parameter $r$ for the same purpose. More fundamentally, LoRA and OFT embody two different principles for parameter efficiency: while LoRA leverages low-rankness to decrease the number of tunable parameters, OFT achieves efficiency via structured sparsity, specifically through block-diagonal orthogonality.

\textbf{Why OFT converges faster}. As indicated in Figure~\ref{fig:toy_ae}, the most impactful way to alter the semantics involves changing the orientation of neurons—a task OFT is explicitly constructed to handle. Additionally, OFT can be interpreted as finetuning weights while constraining them to a smooth hypersphere manifold, which leads to a more favorable optimization landscape. This observation is supported empirically by \cite{liu2021orthogonal}.

\textbf{Why not minimize hyperspherical energy}. Unlike the approach of \cite{liu2021orthogonal}, we do not target minimization of hyperspherical energy. For classification tasks, non-redundant neurons are preferable. Achieving minimal hyperspherical energy enforces uniform distribution of neurons over the hypersphere, which is not aligned with the goals of finetuning, as such uniformity could undermine the pretrained features.

\textbf{Balancing flexibility and regularity in finetuning.} Our investigation reveals an intrinsic trade-off between flexibility and regularity in finetuning strategies. While standard finetuning offers maximal flexibility, it often results in unstable outcomes and heightened risk of model collapse. In contrast, the straightforward approach of OFT achieves an effective compromise by maintaining the angles between neurons. The block-diagonal parameterization can be interpreted as a more rigorous form of regularization for the orthogonal transformation matrix.

\textbf{Inference efficiency maintained.} Distinct from approaches such as ControlNet, the OFT framework does not impose extra inference costs on the finetuned model. During inference, it suffices to left-multiply the learned orthogonal matrix $\bm{R}$ with the pretrained weights $\bm{W}^0$ to directly produce a new effective weight matrix $\thickmuskip=2mu \medmuskip=2mu \bm{W}=\bm{R}\bm{W}^0$. As a result, the inference performance is identical to that of the original pretrained model.

\section{Experiments and Results}

\textbf{General experimental setup.} Our experiments utilize Stable Diffusion v1.5~\cite{rombach2022high} as the foundational text-to-image generator. For unbiased assessment, samples from each method are selected randomly. Subject-driven experiments follow the general protocol outlined by DreamBooth~\cite{ruiz2023dreambooth}, while approaches for controllable generation align with ControlNet~\cite{zhang2023adding} and T2I-Adapter~\cite{mou2023t2i}. To directly compare with LoRA, we restrict the application of OFT or COFT to the same network layers targeted by LoRA. Further implementation specifics and additional experimental outcomes can be found in Appendix~\ref{app:settings}.

\subsection{Subject-driven Generation}

\textbf{Experimental configuration.} We benchmark against DreamBooth~\cite{ruiz2023dreambooth} and LoRA~\cite{hulora2022}. Every method employs the same loss formulation utilized in DreamBooth. Adhering to the protocols from the respective original works, both DreamBooth and LoRA are run with their optimal hyperparameters. Supplementary results and details appear in Appendix~\ref{app:settings},\ref{app:coft_oft},\ref{app:more_qualitative},\ref{app:failure}.

\textbf{Finetuning stability and convergence}. We begin by assessing the stability and convergence rate of finetuning using DreamBooth, LoRA, OFT, and COFT. Figure~\ref{fig:teaser} and Figure~\ref{fig:db_convergence} display the results of these evaluations. Notably, COFT and OFT demonstrate robust and stable finetuning behavior when applied to the diffusion model. After 400 iterations, both DreamBooth and OFT-based approaches provide effective control, but LoRA exhibits shortcomings in maintaining the subject's identity. At 2000 iterations, DreamBooth starts producing degraded, collapsed outputs, and LoRA fails to render the subject in a yellow shirt—mistakenly introducing yellow fur instead. In contrast, OFT and COFT continue to provide stable and coherent control over image generation even at this extended training length. These observations attest to the rapid yet dependable convergence properties of our OFT approach for subject-driven image synthesis. The apparent resilience of OFT and COFT to the specific choice of iteration number is particularly advantageous, as it reduces the need for extensive hyperparameter tuning across different subjects. This allows for simply selecting a sufficiently large number of iterations for both OFT and COFT without risk, bypassing intricate adjustment procedures. When COFT is used with an appropriately chosen $\epsilon$, the selection of both learning rate and total iterations becomes straightforward and requires minimal manual intervention.

\textbf{Quantitative comparison}. In accordance with \cite{ruiz2023dreambooth}, we perform a quantitative assessment that includes subject fidelity (measured by DINO~\cite{caron2021emerging} and CLIP-I~\cite{radford2021learning}), adherence to text prompts (CLIP-T~\cite{radford2021learning}), and diversity among samples (LPIPS~\cite{zhang2018unreasonable}). For CLIP-I, we compute the mean pairwise cosine similarity between the CLIP embeddings of generated and original images, whereas DINO uses ViT S/16 DINO embeddings for a similar calculation. CLIP-T is assessed by measuring the average cosine similarity between CLIP embeddings of the input text prompt and the resulting images. Additionally, LPIPS is used to evaluate the average cosine similarity among generated images corresponding to the same subject and prompt. Results summarized in Table~\ref{table:dreambooth_final} indicate that both COFT and OFT yield substantial improvements over DreamBooth and LoRA in terms of DINO and CLIP-I, while maintaining similar or marginally better performance in terms of text prompt fidelity and diversity. To ensure reliable results, we finetune each model configuration 30 times using different random seeds, minimizing variability due to chance. Overall, these outcomes confirm that the OFT approach not only excels in achieving fast and stable convergence, but also consistently delivers superior final performance metrics.

\textbf{Qualitative comparison}. To provide a clearer picture of OFT's advantages, we present randomly selected subject-driven generation samples in Figure~\ref{fig:dreambooth_final}. All samples originate from the same finetuned base model for each technique to ensure a fair comparison; thus, individual text prompts are not separately tuned for specific outcomes. For each method, we report results using the model that obtained the highest validation CLIP metrics. From the qualitative results in Figure~\ref{fig:dreambooth_final}, it is evident that both OFT and COFT maintain strong semantic fidelity to the subject, whereas LoRA frequently fails to preserve the subject’s identity (\emph{e.g.}, the subject is completely unrecognizable in the bowl example with LoRA). Additionally, OFT and COFT offer substantially better prompt-based control, while DreamBooth, though effective in retaining subject identity, often does not accurately generate images corresponding to the prompt (\emph{e.g.}, first row of the bowl example). These comparisons reveal that our OFT framework excels in both subject preservation and controllability. Furthermore, OFT demonstrates robustness to the number of iterations—showing stable performance even when the iteration count is high—a property not shared by DreamBooth or LoRA. Additional qualitative results can be found in Appendix~\ref{app:more_qualitative}. Appendix~\ref{app:human} details a human study further supporting OFT’s superiority.

\subsection{Controllable Generation}

\textbf{Settings}. Our baseline comparisons include ControlNet~\cite{zhang2023adding}, T2I-Adapter~\cite{mou2023t2i}, and LoRA~\cite{hulora2022}. The main paper addresses three complex controllable generation challenges: canny edge-to-image (C2I) using the COCO dataset~\cite{lin2014microsoft}, segmentation map-to-image (S2I) on ADE20K~\cite{zhou2017scene}, and landmark-to-face (L2F) on CelebA-HQ~\cite{karras2018progressive,xia2021tedigan}. Each approach is used to finetune Stable Diffusion~(SD) v1.5 for 20 epochs on these datasets. Further qualitative examples are presented in Appendix~\ref{app:more_qualitative},\ref{app:more_control_task},\ref{app:failure}.

\textbf{Convergence}. To assess the convergence properties of ControlNet, T2I-Adapter, LoRA, and COFT on the L2F task, we conduct both quantitative and qualitative analyses. For our quantitative assessment, we employ the mean $\ell_2$ distance between predicted and control face landmarks as the metric. Figure~\ref{fig:face_convergence} illustrates the trajectory of face landmark errors as the models are fine-tuned over increasing epochs. Notably, COFT and OFT display a markedly accelerated convergence rate in comparison to the alternatives. For example, LoRA requires 20 epochs to reach the same level of performance that our OFT framework achieves in only 8 epochs. It is important to highlight that both OFT and COFT possess a similar parameter count to LoRA (and far fewer than ControlNet), yet surpass other approaches in convergence efficiency. Furthermore, the rapid convergence observed with OFT is corroborated by the examples in Figure~\ref{fig:teaser}; in the rightmost example, OFT demonstrates substantially greater data efficiency than both ControlNet and LoRA, achieving satisfactory convergence with merely 5\% of the ADE20K dataset. 

For our qualitative study, we focus on OFT, COFT, and ControlNet, as ControlNet’s landmark error most closely matches our results. The visual evidence provided in Figure~\ref{fig:face_convergence_qual} indicates that OFT and COFT both demonstrate consistent and stable convergence, progressively adjusting the generated face poses to align with the control landmarks. In contrast, ControlNet exhibits a delayed controllability, with a pronounced improvement only manifesting after the 8-th epoch—this abrupt transition aligns with the convergence dynamics discussed in \cite{zhang2023adding}. The smoother and more reliable convergence behavior of our OFT framework greatly facilitates practical use, as its training process is predictable and less susceptible to erratic shifts, thereby simplifying the selection of effective hyperparameters.

\textbf{Quantitative comparison}. To assess controllable generation, we propose a control consistency metric, which quantitatively measures how well the generated image reflects the original control input by extracting the corresponding control signal from the output and comparing it to the original. Specifically, for the C2I task, IoU and F1 scores are computed; for the S2I task, mean IoU as well as mean and overall accuracy are reported; for the L2F task, we evaluate the mean $\ell_2$ distance between the predicted and input control landmarks. Detailed descriptions of these metrics are presented in Appendix~\ref{app:settings}. All baselines are evaluated using their optimal hyperparameter configurations. As indicated in Table~\ref{table:control_gen}, both OFT and COFT demonstrate substantially higher accuracy and robustness in control compared to other approaches. Adapter-based baselines (\emph{e.g.}, T2I-Adapter and ControlNet) tend to require more iterations to converge and ultimately deliver inferior performance. Relative to ControlNet, LoRA achieves superior outcomes in the S2I setting but underperforms in C2I and L2F contexts. More broadly, even after thorough hyperparameter optimization, LoRA’s maximum attainable performance remains limited. By contrast, our OFT framework continues to benefit as its parameter count increases, and has yet to reach its performance plateau. Notably, our quantitative assessment framework for controllable generation represents a novel contribution, enabling precise measurement of a fine-tuned model’s control ability (contingent on the accuracy of the pre-existing segmentation/detection network).

\textbf{Qualitative comparison}. We further conduct a qualitative analysis comparing OFT and COFT with leading contemporary approaches such as ControlNet, T2I-Adapter, and LoRA. As illustrated by the randomly sampled results in Figure~\ref{fig:control_final}, both OFT and COFT consistently produce images that are not only photorealistic and high in fidelity, but also conform precisely to the specified control signals. On the S2I task, for instance, LoRA is unable to faithfully render images based on the provided segmentation maps, whereas ControlNet, OFT, and COFT exhibit robust control over output generation. Notably, when juxtaposed with ControlNet, OFT and COFT are capable of synthesizing high-fidelity outputs characterized by richer details and more coherent geometric arrangements, all while employing significantly fewer parameters. In the C2I task, OFT and COFT successfully generate semantically coherent images from crude Canny edge inputs, in contrast to the inferior performance of T2I-Adapter and LoRA. When addressing the L2F task, our proposed approach delivers the highest accuracy in pose guidance for generated facial images, even in the context of complex facial orientations. Across these three control scenarios, results consistently demonstrate that OFT and COFT produce images of superior qualitative quality relative to other state-of-the-art baselines, underscoring the robustness of our OFT framework for controllable image synthesis. For broader qualitative insights, additional examples across all three control tasks can be found in Appendix~\ref{app:control_exp}. Furthermore, Appendix~\ref{app:more_control_task} presents evidence of OFT's versatility in a broader spectrum of control scenarios, including dense pose to human body, sketch to image, and depth to image translation.

\section{Concluding Remarks and Open Problems}

Inspired by the finding that the angular relationships between neurons play a pivotal role in shaping visual semantics, we introduce a straightforward yet powerful finetuning strategy—orthogonal finetuning—for exerting control over text-to-image diffusion models. Our approach is specifically designed for two key applications: subject-driven generation and controllable generation. In comparison with current approaches, OFT achieves superior controllability and stable finetuning while utilizing fewer parameters. Crucially, OFT maintains the same inference cost as baseline models, facilitating efficient deployment.

OFT gives rise to several intriguing avenues for further investigation. To enforce orthogonality, OFT employs Cayley parametrization, which necessitates a matrix inversion step. This requirement poses certain constraints on the scalability of the method. Although our use of block diagonal parametrization alleviates this issue to some extent, devising faster, differentiable matrix inversion techniques remains an open challenge. Additionally, OFT holds promise for compositionality: the orthogonal matrices generated through distinct OFT finetuning procedures can be multiplied together, yielding another orthogonal matrix. Exploring whether this multiplicative set of orthogonal matrices retains the specific capabilities imparted by each individual downstream task is an open research question. Lastly, OFT's parameter efficiency is inherently tied to the block diagonal structure, which can introduce additional biases and curtail flexibility. Developing alternative strategies to further enhance parameter efficiency without sacrificing flexibility or incurring bias represents an important problem for future work.